% !TeX program = XeLaTeX
\documentclass[12pt,a4paper,UTF8]{ctexart}

\usepackage[a4paper,margin=2.2cm]{geometry}
\usepackage{amsmath,amssymb,amsfonts}
\usepackage{booktabs,longtable,array,multirow}
\usepackage{graphicx}
\usepackage{hyperref}
\usepackage{xcolor}
\usepackage{enumitem}
\usepackage{listings}
\usepackage{setspace}
\usepackage{float}
\usepackage{fontspec}

\IfFontExistsTF{Noto Serif CJK TC}{
    \setCJKmainfont{Noto Serif CJK TC}
    \setCJKsansfont{Noto Sans CJK TC}
}{
    \IfFontExistsTF{Noto Serif CJK SC}{
        \setCJKmainfont{Noto Serif CJK SC}
        \setCJKsansfont{Noto Sans CJK SC}
    }{
        \setCJKmainfont{FandolSong-Regular}
        \setCJKsansfont{FandolHei-Regular}
    }
}

\setmainfont{TeX Gyre Termes}
\setsansfont{TeX Gyre Heros}
\setmonofont{Latin Modern Mono}

\hypersetup{
    colorlinks=true,
    linkcolor=blue!60!black,
    urlcolor=blue!60!black,
    citecolor=blue!60!black
}

\lstdefinestyle{pythonstyle}{
    language=Python,
    basicstyle=\ttfamily\small,
    keywordstyle=\color{blue!70!black},
    commentstyle=\color{green!40!black},
    stringstyle=\color{red!60!black},
    showstringspaces=false,
    breaklines=true,
    frame=single,
    columns=fullflexible
}

\title{RedRhex 彈簧腿節能導向強化學習 Reward 整合報告}
\author{Codex Research Assistant}
\date{2026-03-31}

\begin{document}
\maketitle
\tableofcontents
\newpage

\section{研究問題重述}

\subsection{系統背景}

本研究對象為 RedRhex 六足 wheg 機器人。每條腿包含三種功能不同的關節：

\begin{enumerate}[leftmargin=2em]
    \item \textbf{main drive 關節}：速度控制，負責 wheg 旋轉與主要推進。
    \item \textbf{ABAD 關節}：位置控制，負責外展/內收與側向落點調整。
    \item \textbf{damper/spring-like 關節}：被動關節，負責吸震、儲能與釋能。
\end{enumerate}

目前使用 Isaac Lab 環境：
\begin{itemize}[leftmargin=2em]
    \item \texttt{Template-Redrhex-Direct-v0}
    \item \texttt{Template-Redrhex-ForwardFast-Direct-v0}
\end{itemize}

主要目標不是只讓策略學會「能走」，而是讓策略學會：
\begin{enumerate}[leftmargin=2em]
    \item 追蹤命令速度；
    \item 保持穩定與不摔倒；
    \item 避免暴力輸出與 reward hacking；
    \item 降低有源能耗；
    \item 實際利用彈簧腿儲能與釋能優勢。
\end{enumerate}

\subsection{核心研究問題}

本研究要回答的核心問題為：

\begin{quote}
如何把 RedRhex 的被動彈簧腿機制，正式整合進 Isaac Lab 強化學習 reward，使策略不僅能 locomote，還能在保持穩定與追蹤性能的前提下，學會更低能耗且更具彈簧回能特性的步態？
\end{quote}

\section{物理與數學推導}

\subsection{關節角度、角速度與力矩}

對每條腿定義：
\begin{align}
q_m &: \text{main drive 關節角}, \\
q_a &: \text{ABAD 關節角}, \\
q_s &: \text{spring/damper 關節角}.
\end{align}

角速度定義為
\begin{equation}
\omega = \dot{q} = \frac{dq}{dt}.
\end{equation}

\subsection{Main drive 力矩模型}

對 main drive 關節，Isaac Lab implicit actuator 可近似為速度控制器：
\begin{equation}
\tau_{m,i} \approx b_m \left( \omega^{cmd}_{m,i} - \omega_{m,i} \right),
\end{equation}
其中本專案採用
\begin{equation}
b_m = 50, \qquad |\tau_{m,i}| \le 100 \text{ N$\cdot$m}.
\end{equation}

因此 fallback 力矩估計為
\begin{equation}
\tau^{est}_{m,i} =
\mathrm{clip}\left(
50(\omega^{cmd}_{m,i} - \omega_{m,i}), -100, 100
\right).
\end{equation}

\subsection{ABAD 力矩模型}

對 ABAD 關節，採位置控制近似：
\begin{equation}
\tau_{a,i} \approx k_a(q^{cmd}_{a,i} - q_{a,i}) - b_a \omega_{a,i},
\end{equation}
其中
\begin{equation}
k_a = 40, \qquad b_a = 4, \qquad |\tau_{a,i}| \le 8 \text{ N$\cdot$m}.
\end{equation}

因此 fallback 估計為
\begin{equation}
\tau^{est}_{a,i} =
\mathrm{clip}\left(
40(q^{cmd}_{a,i} - q_{a,i}) - 4\omega_{a,i}, -8, 8
\right).
\end{equation}

\subsection{被動彈簧與阻尼模型}

令 damper 偏轉為
\begin{equation}
\Delta q_{s,i} = q_{s,i} - q^{rest}_{s,i}.
\end{equation}

則彈簧力矩與阻尼力矩可寫為
\begin{align}
\tau^{spring}_{s,i} &= -k_s \Delta q_{s,i}, \\
\tau^{damp}_{s,i} &= -d_s \omega_{s,i},
\end{align}
其中目前設計使用
\begin{equation}
k_s = 200 \text{ N$\cdot$m/rad}, \qquad d_s = 20 \text{ N$\cdot$m$\cdot$s/rad}.
\end{equation}

\subsection{機械功率與機械能}

單一關節的瞬時機械功率為
\begin{equation}
P_i = \tau_i \omega_i.
\end{equation}

由於本研究希望估計有源輸出成本，而非淨做功，因此使用絕對值：
\begin{align}
P_{main} &= \sum_{i \in main} |\tau_i \omega_i|, \\
P_{abad} &= \sum_{j \in abad} |\tau_j \omega_j|, \\
P_{act} &= P_{main} + P_{abad}.
\end{align}

每一步機械能 proxy 為
\begin{equation}
E_{step} = P_{act}\Delta t,
\end{equation}
單回合則為
\begin{equation}
E_{episode} = \sum_t P_{act}(t)\Delta t.
\end{equation}

\subsection{電流與電功率估計}

若真機可量到 motor current，則可由
\begin{equation}
\tau_{joint} = N \eta_g K_t I
\end{equation}
得
\begin{equation}
I = \frac{\tau_{joint}}{N \eta_g K_t}.
\end{equation}

電功率可近似為
\begin{equation}
P_{elec} \approx I^2R + \frac{\tau \omega}{\eta_m \eta_d}.
\end{equation}

因此即使不知道完整電機常數，仍可用兩個重要 proxy：
\begin{enumerate}[leftmargin=2em]
    \item $|\tau \omega|$：機械功率 proxy；
    \item $\tau^2$：銅損 proxy。
\end{enumerate}

\subsection{彈簧位能、釋能與耗散}

彈簧位能為
\begin{equation}
E_{s,i} = \frac{1}{2}k_s(\Delta q_{s,i})^2.
\end{equation}

位能變化率為
\begin{equation}
\dot{E}_{s,i} = k_s \Delta q_{s,i} \dot{q}_{s,i}.
\end{equation}

定義儲能與釋能功率為
\begin{align}
P_{store} &= \sum_i \max(\dot{E}_{s,i}, 0), \\
P_{release} &= \sum_i \max(-\dot{E}_{s,i}, 0).
\end{align}

阻尼耗散功率為
\begin{equation}
P_{diss} = d_s \sum_i \omega_{s,i}^2.
\end{equation}

\subsection{等效運動速度與 CoT proxy}

本研究為了讓 yaw 模式也能合理做能耗比較，引入等效運動速度
\begin{equation}
v_{eq} = \sqrt{v_x^2 + v_y^2 + (r_{yaw}\omega_z)^2},
\end{equation}
其中
\begin{equation}
r_{yaw} = 0.18 \text{ m}.
\end{equation}

定義 CoT proxy 為
\begin{equation}
CoT^* = \frac{P_{act}}{mg(v_{eq} + \epsilon)}.
\end{equation}

\subsection{彈簧回收率與活用度}

本研究實際使用的 spring recovery ratio 為
\begin{equation}
\eta_{rec}
=
\mathrm{clamp}\left(
\frac{P_{release}}{P_{main} + \alpha P_{abad} + \epsilon_{rec}},
0,1
\right),
\end{equation}
其中
\begin{equation}
\alpha = 0.35.
\end{equation}

彈簧活用度則採用六條腿偏轉的標準差：
\begin{equation}
U_s = \mathrm{std}(\Delta q_{s,1}, \ldots, \Delta q_{s,6}).
\end{equation}

\section{模擬 / 真機感測映射}

\subsection{模擬中可直接取得的量}

\begin{longtable}{p{3cm}p{4cm}p{6cm}}
\toprule
物理量 & 模擬中理想讀取方式 & 備註 \\
\midrule
關節角度 $q$ & \texttt{self.joint\_pos[:, idx]} & 直接由 articulation data 提供 \\
關節角速度 $\omega$ & \texttt{self.joint\_vel[:, idx]} & 直接可用 \\
力矩 $\tau$ & \texttt{applied\_torque} / \texttt{computed\_torque} / \texttt{joint\_torque} & 視 Isaac Lab 版本而定 \\
base 線速度 & \texttt{self.base\_lin\_vel} & 可直接用於 tracking 與 CoT \\
base 角速度 & \texttt{self.base\_ang\_vel} & yaw 模式很重要 \\
damper rest pose & \texttt{cfg.robot\_cfg.init\_state.joint\_pos} & 用於計算 $\Delta q_s$ \\
接觸資訊 & \texttt{\_current\_leg\_in\_stance} & 目前為 phase-based proxy \\
\bottomrule
\end{longtable}

\subsection{真機感測的三層方案}

\subsubsection*{Ideal sensing}
\begin{itemize}[leftmargin=2em]
    \item encoder 量測 $q$；
    \item torque transducer 量測 $\tau$；
    \item motion capture 或 state estimator 量測 $(v_x,v_y,\omega_z)$；
    \item foot force plate 或 load cell 量測 GRF/contact。
\end{itemize}

\subsubsection*{Practical sensing}
\begin{itemize}[leftmargin=2em]
    \item encoder 差分估 $\omega$；
    \item motor current 轉 torque；
    \item IMU + odometry fusion 估 base velocity；
    \item 利用 linkage 幾何回推 damper 偏轉。
\end{itemize}

\subsubsection*{Fallback proxy}
\begin{itemize}[leftmargin=2em]
    \item 若無 torque sensor 與 current，退回 controller-based torque estimate；
    \item 若無真實 contact sensor，退回 gait phase/contact proxy；
    \item 若無彈簧位移量測，退回相對幾何壓縮 proxy。
\end{itemize}

\section{Reward 設計}

\subsection{設計哲學}

本設計遵循以下優先順序：
\begin{equation}
\text{stability} > \text{tracking} > \text{anti-cheat} > \text{energy} > \text{spring use}.
\end{equation}

也就是說：
\begin{itemize}[leftmargin=2em]
    \item 先保命；
    \item 再追蹤；
    \item 最後才對能效與彈簧利用做二階優化。
\end{itemize}

\subsection{E1: Power Efficiency}

\begin{equation}
r_{E1}
=
-\tanh\left(
\frac{P_{act}}{(v_{eq}+\epsilon_v)s_{tanh}}
\right)
\cdot w_{power}
\cdot \mathbb{1}_{healthy}
\end{equation}

本專案目前設定：
\begin{align}
w_{power} &= 0.3, \\
\epsilon_v &= 0.1, \\
s_{tanh} &= 500.0.
\end{align}

\subsection{E2: Spring Recovery}

\begin{equation}
r_{E2}
=
\eta_{rec}\cdot w_{rec}
\cdot \mathbb{1}_{healthy}
\cdot \mathbb{1}_{cmd}
\cdot g_{track}
\end{equation}

其中：
\begin{itemize}[leftmargin=2em]
    \item $\mathbb{1}_{cmd}$：命令等效速度超過門檻；
    \item $g_{track}$：實際等效速度相對命令等效速度的達成比例；
    \item 並且僅在 stance contact proxy 成立的腿上計入 $P_{release}$ 與 $P_{store}$。
\end{itemize}

\subsection{E3: Spring Utilization}

\begin{equation}
r_{E3}
=
\mathrm{clamp}\left(
\frac{U_s}{\Delta q_{max}}, 0, 1
\right)
\cdot w_{util}
\cdot \mathbb{1}_{healthy}
\cdot \mathbb{1}_{cmd}
\cdot g_{track}
\end{equation}

其中
\begin{equation}
\Delta q_{max} = 0.3 \text{ rad}.
\end{equation}

\subsection{E4: Torque Penalty}

\begin{equation}
r_{E4}
=
\left(
\sum \tau_m^2 + \beta \sum \tau_a^2
\right) w_{\tau}
\end{equation}
其中
\begin{equation}
\beta = 0.5, \qquad w_{\tau} = -10^{-4}.
\end{equation}

\subsection{為何需要 spring reward gating}

若只獎勵 spring release 或 spring utilization，策略可能學會：
\begin{itemize}[leftmargin=2em]
    \item 原地抖動 damper；
    \item 空中甩腿刷彈簧偏轉；
    \item 不追蹤命令，卻拿到正面 spring reward。
\end{itemize}

因此目前正式實作採用
\begin{equation}
\texttt{spring\_reward\_gate}
=
\texttt{healthy\_gate}
\cdot
\texttt{cmd\_motion\_gate}
\cdot
\texttt{motion\_tracking\_gate}.
\end{equation}

\section{模擬實作與程式修改}

\subsection{\texttt{redrhex\_env\_cfg.py} 中新增的參數}

\begin{lstlisting}[style=pythonstyle,caption={節能 reward 權重與物理常數}]
"power_efficiency": 0.3,
"power_efficiency_eps": 0.1,
"power_efficiency_tanh_scale": 500.0,
"spring_recovery": 0.4,
"spring_recovery_eps": 0.01,
"spring_recovery_abad_weight": 0.35,
"spring_utilization": 0.2,
"spring_util_max_deflection": 0.3,
"torque_penalty": -0.0001,
"torque_penalty_abad_weight": 0.5,

damper_stiffness = 200.0
damper_damping = 20.0
robot_mass_kg = 14.0
energy_velocity_yaw_radius = 0.18
energy_min_command_motion = 0.05
main_drive_torque_estimate_damping = 50.0
main_drive_torque_estimate_limit = 100.0
abad_torque_estimate_stiffness = 40.0
abad_torque_estimate_damping = 4.0
abad_torque_estimate_limit = 8.0
\end{lstlisting}

\subsection{\texttt{redrhex\_env.py} 中新增的 helper}

\begin{lstlisting}[style=pythonstyle,caption={等效運動速度 helper}]
def _compute_energy_equivalent_speed(self, lin_xy, yaw_rate):
    yaw_equiv = self._energy_yaw_radius * yaw_rate
    return torch.sqrt(torch.sum(torch.square(lin_xy), dim=1) + torch.square(yaw_equiv))
\end{lstlisting}

\begin{lstlisting}[style=pythonstyle,caption={主動關節力矩讀取與 fallback}]
def _get_active_joint_torques(self, main_drive_vel, abad_vel):
    for attr_name in ("applied_torque", "computed_torque", "joint_torque"):
        candidate = getattr(self.robot.data, attr_name, None)
        if isinstance(candidate, torch.Tensor) and candidate.ndim == 2:
            return candidate[:, self._main_drive_indices], candidate[:, self._abad_indices]

    main_torques = self._main_drive_torque_estimate_damping * (self._target_drive_vel - main_drive_vel)
    main_torques = torch.clamp(main_torques, -self._main_drive_torque_estimate_limit, self._main_drive_torque_estimate_limit)

    abad_pos_err = self._target_abad_pos - self.joint_pos[:, self._abad_indices]
    abad_torques = self._abad_torque_estimate_stiffness * abad_pos_err - self._abad_torque_estimate_damping * abad_vel
    abad_torques = torch.clamp(abad_torques, -self._abad_torque_estimate_limit, self._abad_torque_estimate_limit)
    return main_torques, abad_torques
\end{lstlisting}

\subsection{ABAD 目標位置 cache}

\begin{lstlisting}[style=pythonstyle,caption={ABAD target cache}]
target_abad_pos = torch.clamp(target_abad_pos, min=-abad_limit, max=abad_limit)
self._target_abad_pos = target_abad_pos.clone()
self.robot.set_joint_position_target(target_abad_pos, joint_ids=self._abad_indices)
\end{lstlisting}

\subsection{reward 核心節能片段}

\begin{lstlisting}[style=pythonstyle,caption={節能 reward 核心計算}]
main_torques, abad_torques = self._get_active_joint_torques(main_drive_vel, abad_vel)

mech_power_main = torch.sum(torch.abs(main_torques * main_drive_vel), dim=1)
mech_power_abad = torch.sum(torch.abs(abad_torques * abad_vel), dim=1)
total_mech_power = mech_power_main + mech_power_abad

actual_motion_speed = self._compute_energy_equivalent_speed(actual_lin, actual_wz)
cmd_motion_speed = self._compute_energy_equivalent_speed(cmd_lin, cmd_wz)

rew_power_efficiency = -torch.tanh(
    (total_mech_power / (actual_motion_speed + power_eff_eps)) / power_eff_tanh_scale
) * scales.get("power_efficiency", 0.3)

if hasattr(self, "_current_leg_in_stance") and self._current_leg_in_stance.shape == damper_pos.shape:
    spring_contact_mask = self._current_leg_in_stance.float()
else:
    spring_contact_mask = torch.ones_like(damper_pos)

spring_release = torch.sum(torch.clamp(-spring_power, min=0.0) * spring_contact_mask, dim=1)
spring_store = torch.sum(torch.clamp(spring_power, min=0.0) * spring_contact_mask, dim=1)

cmd_motion_gate = (cmd_motion_speed > self._energy_min_cmd_motion).float()
motion_tracking_gate = torch.clamp(
    actual_motion_speed / torch.clamp(cmd_motion_speed, min=power_eff_eps),
    min=0.0,
    max=1.0,
)
spring_reward_gate = healthy_gate * cmd_motion_gate * motion_tracking_gate
\end{lstlisting}

\subsection{\texttt{eval\_command\_sweep.py} 中新增的 energy evaluation}

\begin{lstlisting}[style=pythonstyle,caption={energy KPI 收集 helper}]
def collect_energy_metrics(unwrapped_env, actual_vx, actual_vy, actual_wz):
    lin_xy = torch.stack((actual_vx, actual_vy), dim=1)
    yaw_radius = float(getattr(unwrapped_env, "_energy_yaw_radius", 0.18))
    motion_speed = torch.sqrt(torch.sum(torch.square(lin_xy), dim=1) + torch.square(yaw_radius * actual_wz))
    return {
        "motion_speed": motion_speed,
        "mech_power_main": main_power,
        "mech_power_abad": abad_power,
        "mech_power_total": total_power,
        "cot_proxy": cot_proxy,
        "spring_energy": spring_energy,
        "spring_release": spring_release,
        "spring_store": spring_store,
        "spring_recovery_ratio": spring_recovery_ratio,
        "damper_dissipation": damper_dissipation,
    }
\end{lstlisting}

每個 command 目前都會記錄：
\begin{itemize}[leftmargin=2em]
    \item \texttt{energy\_mech\_power\_main\_mean}
    \item \texttt{energy\_mech\_power\_total\_mean}
    \item \texttt{energy\_cost\_of\_transport\_proxy}
    \item \texttt{energy\_spring\_energy\_mean}
    \item \texttt{energy\_spring\_release\_power\_mean}
    \item \texttt{energy\_spring\_store\_power\_mean}
    \item \texttt{energy\_spring\_recovery\_ratio}
    \item \texttt{energy\_motion\_speed\_mean}
    \item \texttt{energy\_power\_per\_motion}
\end{itemize}

\section{權重設計與調參建議}

\begin{longtable}{p{4cm}p{2cm}p{2cm}p{6cm}}
\toprule
項目 & Default & ForwardFast & 調參原則 \\
\midrule
\texttt{power\_efficiency} & 0.30 & 0.15 & 先保守，tracking 穩後再加強 \\
\texttt{spring\_recovery} & 0.40 & 0.20 & 是最核心的 spring reward \\
\texttt{spring\_utilization} & 0.20 & 0.10 & 小權重，避免刷偏轉 \\
\texttt{torque\_penalty} & -1e-4 & -5e-5 & 只抑制極端有源輸出 \\
\texttt{torque\_penalty\_abad\_weight} & 0.50 & 0.50 & 避免 ABAD 被過度壓制 \\
\bottomrule
\end{longtable}

推薦調整順序：
\begin{enumerate}[leftmargin=2em]
    \item baseline tracking/stability；
    \item 加 torque penalty；
    \item 加 power efficiency；
    \item 最後加 spring recovery / spring utilization。
\end{enumerate}

\section{Ablation 與驗證流程}

\subsection{四版本比較}

\begin{enumerate}[leftmargin=2em]
    \item A：baseline；
    \item B：baseline + torque penalty；
    \item C：baseline + power efficiency；
    \item D：baseline + power efficiency + spring recovery/utilization。
\end{enumerate}

\subsection{TensorBoard 必看曲線}

\begin{itemize}[leftmargin=2em]
    \item \texttt{rew\_power\_efficiency}
    \item \texttt{rew\_spring\_recovery}
    \item \texttt{rew\_spring\_utilization}
    \item \texttt{rew\_torque\_penalty}
    \item \texttt{diag\_mech\_power\_total}
    \item \texttt{diag\_cost\_of\_transport}
    \item \texttt{diag\_spring\_recovery\_ratio}
    \item \texttt{diag\_motion\_speed\_equiv}
    \item \texttt{diag\_cmd\_motion\_speed\_equiv}
\end{itemize}

\subsection{如何判斷「真的變省能」}

以下條件至少應同時成立：
\begin{enumerate}[leftmargin=2em]
    \item tracking quality 不顯著下降；
    \item success ratio 不顯著下降；
    \item fall rate 不惡化；
    \item $CoT^*$ 下降；
    \item power-per-motion 下降。
\end{enumerate}

若只看到功率下降，但 motion speed 也大幅下降，則不能聲稱真正更高效。

\section{風險與注意事項}

\begin{itemize}[leftmargin=2em]
    \item 目前 damper contact gating 仍是 phase-based proxy，不是真實 GRF。
    \item 真機若缺 torque/current sensing，很多電功率結果只能作 proxy。
    \item \texttt{spring\_utilization} 是結構 proxy，不是直接機械能收益。
    \item yaw 等效半徑 $r_{yaw}$ 是工程近似，未來可再做系統辨識。
    \item 本次已完成程式實作與靜態語法檢查，但未在本 shell 內直接做 Isaac Lab runtime rollout。
\end{itemize}

\section{結論}

本研究已將 RedRhex 的節能議題從概念層推進到可執行的 Isaac Lab reward engineering 實作。其核心貢獻包括：
\begin{enumerate}[leftmargin=2em]
    \item 以物理一致的方式建構 main drive、ABAD、damper 的能量模型；
    \item 以 $v_{eq}$ 解決 pure yaw 任務的能耗正規化問題；
    \item 將 spring recovery 與 spring utilization 納入 reward，並加入 anti-hacking gating；
    \item 在 evaluation pipeline 中加入 per-command 與 per-skill energy KPI；
    \item 將所有設計落實到 \texttt{redrhex\_env\_cfg.py}、\texttt{redrhex\_env.py} 與 \texttt{eval\_command\_sweep.py}。
\end{enumerate}

因此，這套方法不僅能回答「機器人有沒有走起來」，也能回答更重要的研究問題：
\begin{quote}
機器人是否在保持 locomotion 任務品質的前提下，真的學會用更低的有源功率，並更有效地利用被動彈簧腿機制？
\end{quote}

\end{document}
